\chapter{Complexity — Snapshot}
%%% generated by the blueprint pipeline from TCSlib.Complexity.CookLevin.Snapshot
%%%%%%%%%%%%%%%%%%%%%%%%%%%%%%%%%%%%%%%%%%%%%%%%%%%%%%%%%%%%%%%%%%%%%%%%%%%%%%%
\section{Overview}

This module develops the snapshots behind the Cook--Levin tableau of Arora--Barak
[AB09, \S 2.3.4]. A snapshot of a multi-tape Turing machine at a given time is the
constant-size record of its state and the symbols under its heads. The module defines
snapshots, the head schedule of a machine on inputs of length $n$ (read off from the
all-zeros input $0^n$), the last-visit function on work-tape cells, and the finite local
functions computing the next state, the symbol left in a cell, and the emitted output bit
from a snapshot. Its main results show that for an oblivious machine computation is local:
the snapshot at each time is determined by the previous snapshot, the input bit at the
scheduled input position, and the snapshots at the last visits of the current work cells.

%%%%%%%%%%%%%%%%%%%%%%%%%%%%%%%%%%%%%%%%%%%%%%%%%%%%%%%%%%%%%%%%%%%%%%%%%%%%%%%
\section{Declarations}

\begin{definition}[Snapshot of a Turing machine]
\lean{Complexity.Snapshot}
\label{Complexity.Snapshot}
\leanok
\uses{Turing.FinTM}
Let $M$ be a multi-tape Turing machine over the binary alphabet with a finite set $Q$ of
states and $k$ work tapes, whose tape cells each hold a bit or the blank symbol $\sqcup$.
A snapshot of $M$ is a triple $(q, a, b)$ consisting of a state $q \in Q \cup \{\bot\}$,
where $\bot$ marks a halted machine, an input symbol $a \in \{0, 1, \sqcup\}$, and a tuple
$b = (b_0, \dots, b_{k-1})$ of work-tape symbols $b_\tau \in \{0, 1, \sqcup\}$, one per work
tape.
\end{definition}

\begin{definition}[Snapshot of a run at time $t$]
\lean{Complexity.snapshotAt}
\label{Complexity.snapshotAt}
\leanok
\uses{Complexity.Snapshot, Turing.Cfg.inputSymbol, Turing.Cfg.workTapeSymbols, Turing.FinTM, Turing.MultiTapeTM.initCfg, Turing.MultiTapeTM.runFrom}
Let $M$ be a binary multi-tape Turing machine with finite state set $Q$ and $k$ work tapes,
let $x \in \{0,1\}^*$ and let $t \in \mathbb{N}$. The snapshot of $M$ on $x$ at time $t$ is
the triple consisting of the state ($\bot$ if halted), the symbol under the input head, and
the $k$ symbols under the work-tape heads, in the configuration reached by running $M$ for
$t$ steps from its initial configuration on $x$. Since a step leaves a halted configuration
unchanged, this is defined at every time $t$.
\end{definition}

\begin{definition}[Scheduled input-head position]
\lean{Complexity.inputPosAt}
\label{Complexity.inputPosAt}
\leanok
\uses{Turing.FinTM, Turing.MultiTapeTM.initCfg, Turing.MultiTapeTM.runFrom}
Let $M$ be a binary multi-tape Turing machine and let $n, t \in \mathbb{N}$. On an input of
length $n$ the input cells are numbered $0, 1, \dots, n+1$, where cell $p$ with
$1 \le p \le n$ holds the $p$-th input bit and cells $0$ and $n+1$ are blank boundary cells.
The scheduled input position $\mathrm{inpos}_M(n, t) \in \{0, \dots, n+1\}$ is the position
of the input head after $t$ steps of $M$ on the all-zeros input $0^n$.
\end{definition}

\begin{definition}[Scheduled work-head position]
\lean{Complexity.workPosAt}
\label{Complexity.workPosAt}
\leanok
\uses{Turing.FinTM, Turing.MultiTapeTM.initCfg, Turing.MultiTapeTM.runFrom}
Let $M$ be a binary multi-tape Turing machine with $k$ work tapes whose cells are indexed by
$\mathbb{Z}$, and let $n, t \in \mathbb{N}$ and $\tau \in \{0, \dots, k-1\}$. The scheduled
work position $\mathrm{wpos}_M(n, t, \tau) \in \mathbb{Z}$ is the position of the head of
work tape $\tau$ after $t$ steps of $M$ on the all-zeros input $0^n$.
\end{definition}

\begin{definition}[Last previous visit to the current cell]
\lean{Complexity.prevVisit}
\label{Complexity.prevVisit}
\leanok
\uses{Complexity.workPosAt, Turing.FinTM}
Let $M$ be a binary multi-tape Turing machine with $k$ work tapes indexed by $\mathbb{Z}$,
let $n, t \in \mathbb{N}$ and $\tau \in \{0, \dots, k-1\}$, and for $s \in \mathbb{N}$ let
$w(s)$ be the position of the head of work tape $\tau$ after $s$ steps of $M$ on the
all-zeros input $0^n$. The last previous visit $\mathrm{prev}_M(n, t, \tau)$ is the largest
$s < t$ with $w(s) = w(t)$, and is undefined when no such $s$ exists.
\end{definition}

\begin{definition}[Successor state from a snapshot]
\lean{Complexity.stepState}
\label{Complexity.stepState}
\leanok
\uses{Complexity.Snapshot, Turing.FinTM}
Let $M$ be a binary multi-tape Turing machine with finite state set $Q$, $k$ work tapes and
transition function $\delta$, and let $(q, a, b)$ be a snapshot of $M$, consisting of a state
$q \in Q \cup \{\bot\}$ ($\bot$ meaning halted), an input symbol $a \in \{0, 1, \sqcup\}$ and
work-tape symbols $b = (b_0, \dots, b_{k-1})$ in $\{0, 1, \sqcup\}$. Its successor state is
$\bot$ if $q = \bot$, and otherwise the successor state in $Q \cup \{\bot\}$ prescribed by
the transition $\delta(q, a, b)$.
\end{definition}

\begin{definition}[Symbol left in a work cell by a step]
\lean{Complexity.writtenOrKept}
\label{Complexity.writtenOrKept}
\leanok
\uses{Complexity.Snapshot, Turing.FinTM}
Let $M$ be a binary multi-tape Turing machine with finite state set $Q$, $k$ work tapes and
transition function $\delta$, let $(q, a, b)$ be a snapshot of $M$ (a state
$q \in Q \cup \{\bot\}$ with $\bot$ meaning halted, an input symbol $a$, and work-tape symbols
$b_0, \dots, b_{k-1}$, all in $\{0, 1, \sqcup\}$), and let $\tau \in \{0, \dots, k-1\}$. The
symbol left under work head $\tau$ is $b_\tau$ if $q = \bot$; if $q \in Q$, it is the symbol
in $\{0, 1, \sqcup\}$ that the transition $\delta(q, a, b)$ writes on work tape $\tau$, or
$b_\tau$ when that transition writes nothing on tape $\tau$.
\end{definition}

\begin{definition}[Output bit emitted by a step]
\lean{Complexity.emitted}
\label{Complexity.emitted}
\leanok
\uses{Complexity.Snapshot, Turing.FinTM}
Let $M$ be a binary multi-tape Turing machine with finite state set $Q$, $k$ work tapes,
an append-only output tape and transition function $\delta$, and let $(q, a, b)$ be a
snapshot of $M$ (a state $q \in Q \cup \{\bot\}$ with $\bot$ meaning halted, an input symbol
$a$, and work-tape symbols $b$). The emitted symbol of the snapshot is either a bit or
nothing: it is nothing if $q = \bot$, and otherwise it is the optional output bit that the
transition $\delta(q, a, b)$ appends to the output tape.
\end{definition}

\begin{definition}[Boundary-aware input read]
\lean{Complexity.inputBitAt}
\label{Complexity.inputBitAt}
\leanok
For a binary string $x = x_1 x_2 \cdots x_n$ and a position $p \in \mathbb{N}$, the input
symbol of $x$ at position $p$ is the blank symbol $\sqcup$ if $p = 0$ or $p > n$, and is the
bit $x_p$ if $1 \le p \le n$.
\end{definition}

\begin{theorem}[Oblivious head positions are input-independent]
\lean{Complexity.oblivious_schedule_eq}
\label{Complexity.oblivious_schedule_eq}
\leanok
\uses{Complexity.inputPosAt, Complexity.workPosAt, Turing.FinTM, Turing.FinTM.Oblivious, Turing.MultiTapeTM.initCfg, Turing.MultiTapeTM.runFrom}
\difficulty{1}
Let $M$ be a binary multi-tape Turing machine with $k$ work tapes that is oblivious: for any
two inputs of the same length and every $t \in \mathbb{N}$, the input-head position and all
work-head positions after $t$ steps agree on the two inputs. Then for every input
$x \in \{0,1\}^*$ and every $t \in \mathbb{N}$, the input-head position after $t$ steps of
$M$ on $x$ equals the input-head position after $t$ steps of $M$ on the all-zeros input
$0^{|x|}$, and for every work tape $\tau$ the head position of tape $\tau$ after $t$ steps on
$x$ equals its head position after $t$ steps on $0^{|x|}$.
\end{theorem}

\begin{theorem}[The initial snapshot]
\lean{Complexity.snapshotAt_zero}
\label{Complexity.snapshotAt_zero}
\leanok
\uses{Complexity.inputBitAt, Complexity.snapshotAt, Turing.Cfg.inputSymbol, Turing.FinTM, Turing.FinTM.inputSymbol_at, Turing.MultiTapeTM.initCfg, Turing.MultiTapeTM.runFrom_zero}
\difficulty{3}
Let $M$ be a binary multi-tape Turing machine with initial state $q_0$ and $k$ work tapes,
and let $x \in \{0,1\}^*$. Before any step is taken, $M$ on input $x$ is in state $q_0$, the
symbol under its input head is the first bit $x_1$ of $x$ (or the blank $\sqcup$ if $x$ is
empty), and every one of its $k$ work heads reads the blank $\sqcup$.
\end{theorem}

\begin{theorem}[One-step evolution of the state]
\lean{Complexity.snapshotAt_state_succ}
\label{Complexity.snapshotAt_state_succ}
\leanok
\uses{Complexity.snapshotAt, Complexity.stepState, Turing.Action.apply, Turing.FinTM, Turing.MultiTapeTM.initCfg, Turing.MultiTapeTM.runFrom, Turing.MultiTapeTM.runFrom_succ_eq_step', Turing.MultiTapeTM.step}
\difficulty{3}
Let $M$ be a binary multi-tape Turing machine with finite state set $Q$, $k$ work tapes and
transition function $\delta$, let $x \in \{0,1\}^*$ and $t \in \mathbb{N}$, and let
$q \in Q \cup \{\bot\}$ (with $\bot$ meaning halted), $a$ and $b = (b_0, \dots, b_{k-1})$ be
the state, the symbol under the input head, and the symbols under the work heads after $t$
steps of $M$ on $x$. Then the state of $M$ on $x$ after $t + 1$ steps is $\bot$ if
$q = \bot$, and otherwise is the successor state prescribed by the transition
$\delta(q, a, b)$.
\end{theorem}

\begin{theorem}[Input symbol read at the scheduled position]
\lean{Complexity.snapshotAt_inputSymbol}
\label{Complexity.snapshotAt_inputSymbol}
\leanok
\uses{Complexity.inputBitAt, Complexity.inputPosAt, Complexity.oblivious_schedule_eq, Complexity.snapshotAt, Turing.Cfg.inputSymbol, Turing.FinTM, Turing.FinTM.Oblivious, Turing.FinTM.inputSymbol_at, Turing.MultiTapeTM.initCfg, Turing.MultiTapeTM.runFrom}
\difficulty{4}
Let $M$ be a binary multi-tape Turing machine that is oblivious, meaning that on any two
inputs of the same length its input-head and work-head positions after $t$ steps agree for
every $t$. Let $x = x_1 \cdots x_n \in \{0,1\}^n$ and $t \in \mathbb{N}$, number the input
cells $0, 1, \dots, n+1$ with cell $i \in \{1, \dots, n\}$ holding the $i$-th input bit, and
let $p$ be the position of the input head after $t$ steps of $M$ on the all-zeros input
$0^n$. Then the symbol under the input head after $t$ steps of $M$ on $x$ is the blank
$\sqcup$ if $p = 0$ or $p > n$, and is the bit $x_p$ otherwise.
\end{theorem}

\begin{lemma}[One-step update of a work cell]
\lean{Complexity.workCell_succ}
\label{Complexity.workCell_succ}
\leanok
\uses{Complexity.snapshotAt, Complexity.writtenOrKept, Turing.Action.apply_workTapes, Turing.Cfg.workTapeSymbols, Turing.FinTM, Turing.MultiTapeTM.initCfg, Turing.MultiTapeTM.runFrom, Turing.MultiTapeTM.runFrom_succ_eq_step', Turing.MultiTapeTM.step}
\difficulty{4}
Let $M$ be a binary multi-tape Turing machine with finite state set $Q$, transition function
$\delta$ and $k$ work tapes whose cells are indexed by $\mathbb{Z}$; let $x \in \{0,1\}^*$,
$t \in \mathbb{N}$, $\tau \in \{0, \dots, k-1\}$ and $p \in \mathbb{Z}$. Let $q \in Q \cup
\{\bot\}$ (with $\bot$ meaning halted), $a$ and $b = (b_0, \dots, b_{k-1})$ be the state, the
input-head symbol and the work-head symbols after $t$ steps of $M$ on $x$, and let $h$ be the
position of the head of work tape $\tau$ at that time. Then after $t + 1$ steps, cell $p$ of
work tape $\tau$ has the same content as after $t$ steps if $p \ne h$; if $p = h$, its
content is $b_\tau$ when $q = \bot$, and otherwise is the symbol that $\delta(q, a, b)$ writes
on tape $\tau$, or $b_\tau$ when that transition writes nothing on tape $\tau$.
\end{lemma}

\begin{lemma}[Unvisited cells keep their content]
\lean{Complexity.workCell_eq_of_no_visit}
\label{Complexity.workCell_eq_of_no_visit}
\leanok
\uses{Complexity.workCell_succ, Turing.FinTM, Turing.MultiTapeTM.initCfg, Turing.MultiTapeTM.runFrom}
\difficulty{4}
Let $M$ be a binary multi-tape Turing machine with $k$ work tapes whose cells are indexed by
$\mathbb{Z}$, let $x \in \{0,1\}^*$, $\tau \in \{0, \dots, k-1\}$, $p \in \mathbb{Z}$, and
let $a \le b$ be natural numbers. If for every $r$ with $a \le r < b$ the head of work tape
$\tau$ is not at cell $p$ after $r$ steps of $M$ on $x$, then the content of cell $p$ of
work tape $\tau$ after $b$ steps of $M$ on $x$ equals its content after $a$ steps.
\end{lemma}

\begin{lemma}[No previous visit means a fresh cell]
\lean{Complexity.prevVisit_none_no_visit}
\label{Complexity.prevVisit_none_no_visit}
\leanok
\uses{Complexity.prevVisit, Complexity.workPosAt, Turing.FinTM}
\difficulty{3}
Let $M$ be a binary multi-tape Turing machine with $k$ work tapes indexed by $\mathbb{Z}$,
let $n, t \in \mathbb{N}$ and $\tau \in \{0, \dots, k-1\}$, and for $s \in \mathbb{N}$ let
$w(s)$ be the position of the head of work tape $\tau$ after $s$ steps of $M$ on the
all-zeros input $0^n$, and let $\mathrm{prev}$ be the largest $s < t$ with $w(s) = w(t)$,
undefined when there is no such $s$. If $\mathrm{prev}$ is undefined, then $w(r) \ne w(t)$
for every $r < t$.
\end{lemma}

\begin{lemma}[Characterisation of the last previous visit]
\lean{Complexity.prevVisit_some_last}
\label{Complexity.prevVisit_some_last}
\leanok
\uses{Complexity.prevVisit, Complexity.workPosAt, Turing.FinTM}
\difficulty{3}
Let $M$ be a binary multi-tape Turing machine with $k$ work tapes indexed by $\mathbb{Z}$,
let $n, t \in \mathbb{N}$ and $\tau \in \{0, \dots, k-1\}$, and for $s \in \mathbb{N}$ let
$w(s)$ be the position of the head of work tape $\tau$ after $s$ steps of $M$ on the
all-zeros input $0^n$, and let $\mathrm{prev}$ be the largest $s < t$ with $w(s) = w(t)$,
undefined when there is no such $s$. If $\mathrm{prev}$ is defined and equals $s$, then
$s < t$, $w(s) = w(t)$, and $w(r) \ne w(t)$ for every $r$ with $s < r < t$.
\end{lemma}

\begin{theorem}[Locality of work-tape reads]
\lean{Complexity.snapshotAt_workSymbol}
\label{Complexity.snapshotAt_workSymbol}
\leanok
\uses{Complexity.oblivious_schedule_eq, Complexity.prevVisit, Complexity.prevVisit_none_no_visit, Complexity.prevVisit_some_last, Complexity.snapshotAt, Complexity.workCell_eq_of_no_visit, Complexity.workCell_succ, Complexity.workPosAt, Complexity.writtenOrKept, Turing.FinTM, Turing.FinTM.Oblivious, Turing.MultiTapeTM.initCfg, Turing.MultiTapeTM.runFrom}
\difficulty{5}
Let $M$ be a binary multi-tape Turing machine with finite state set $Q$, transition function
$\delta$ and $k$ work tapes indexed by $\mathbb{Z}$, which is oblivious: on any two inputs of
the same length its input-head and work-head positions after $t$ steps agree for every $t$.
Let $x \in \{0,1\}^n$, $t \in \mathbb{N}$ and $\tau \in \{0, \dots, k-1\}$, and let $w(s)$ be
the position of the head of work tape $\tau$ after $s$ steps of $M$ on $0^n$. If no $s < t$
has $w(s) = w(t)$, then the symbol under head $\tau$ after $t$ steps of $M$ on $x$ is the
blank $\sqcup$. Otherwise let $s$ be the largest such time and let $q \in Q \cup \{\bot\}$
(with $\bot$ meaning halted), $a$ and $b = (b_0, \dots, b_{k-1})$ be the state, the
input-head symbol and the work-head symbols after $s$ steps of $M$ on $x$; then the symbol
under head $\tau$ after $t$ steps of $M$ on $x$ is $b_\tau$ if $q = \bot$, and otherwise is
the symbol that $\delta(q, a, b)$ writes on tape $\tau$, or $b_\tau$ if that transition
writes nothing on tape $\tau$.
\end{theorem}
